При разработке этапа сегментации было принято упрощающее предположение о компоновке страницы. Предполагается, что страница может быть представлена как последовательность полноширинных горизонтальных блоков, расположенных сверху вниз и разделённых промежутками фона (см.~рис.~\ref{fig:segment}). Иными словами, основное содержимое страницы организовано в виде полос, каждая из которых занимает практически всю ширину листа и соответствует одному из двух укрупнённых типов: текстовому или нетекстовому.
Такое предположение хорошо соответствует структуре многих научных и учебных изданий, в том числе архивных химических книг, где текстовые абзацы, заголовки, схемы, таблицы и рисунки обычно располагаются последовательно по вертикали.
Даже если конкретный блок содержит не только текст, а, например, иллюстрацию или химическую схему, он всё равно часто образует отдельную горизонтальную область, отделённую от соседних частей страницы промежутками фона.

Выбор именно такой модели компоновки связан с практическими задачами предварительной обработки. На данном этапе важно не выполнить точную логическую вёрстку документа, а выделить крупные однородные фрагменты страницы для последующей маршрутизации: текстовые блоки передаются в OCR-подсистему, а нетекстовые --- в модули анализа химических структур, формул, таблиц или других специальных объектов. Поэтому в работе используется не универсальная страничная сегментация произвольной сложности, а более простой и интерпретируемый эвристический метод, ориентированный на вертикальное разбиение страницы по областям пустого пространства.
Для выделения полноширинных блоков был реализован эвристический алгоритм сегментации, основанный на поиске горизонтальных промежутков между содержательными областями страницы. Основная идея метода состоит в том, что границы между соседними блоками можно определить по строкам изображения, в которых почти отсутствует полезный контент. Такие области интерпретируются как межблочные пустоты и используются в качестве естественных мест разреза страницы.
Работа алгоритма начинается с загрузки исходного изображения страницы и перевода его в градации серого. Затем строится бинарная маска содержимого, в которой элементы страницы отображаются как передний план, а фон --- как нулевая область. Для этого по умолчанию используется адаптивная бинаризация, более устойчивая к неравномерному освещению и локальным изменениям контраста. При необходимости может применяться и фиксированный порог.
После бинаризации выполняется морфологическое замыкание с вытянутым по горизонтали структурным элементом. Назначение этого этапа состоит в том, чтобы слегка объединить близко расположенные штрихи текста и сделать вертикальный профиль страницы более устойчивым. Без такой операции отдельные буквы и мелкие разрывы могли бы приводить к ложным локальным минимумам и, как следствие, к избыточному числу разрезов.

Далее для каждой строки изображения вычисляется количество пикселей, принадлежащих содержимому. Полученная одномерная последовательность образует вертикальный профиль страницы по оси $Y$, отражающий распределение контента сверху вниз. Для подавления случайных колебаний профиль дополнительно сглаживается скользящим окном. После этого строка считается пустой, если число пикселей содержимого в ней не превышает малого порога, задаваемого как доля от ширины страницы. Тем самым критерий пустоты автоматически масштабируется под размер изображения.
На следующем этапе алгоритм выделяет непрерывные диапазоны пустых строк. Из них сохраняются только достаточно высокие промежутки, поскольку слишком узкие пустоты чаще соответствуют межстрочным интервалам внутри текста, а не реальным границам между крупными блоками. Для каждого допустимого промежутка вычисляется центральная координата, которая используется как точка разреза. Вместе с верхней и нижней границами страницы такие точки формируют набор горизонтальных отсечений, по которым страница разбивается на последовательность полноширинных блоков.
После формирования кандидатов блоков выполняется их простая эвристическая классификация на текстовые и нетекстовые. Для этого анализируется бинарная маска каждого блока. В качестве признаков используются общая плотность пикселей содержимого и количество малых связных компонент. Предполагается, что текстовый блок обычно имеет сравнительно невысокую плотность заполнения, поскольку буквы занимают лишь небольшую часть площади, и одновременно содержит большое число мелких компонент, соответствующих отдельным символам и их частям. Нетекстовые области, напротив, чаще оказываются либо более плотными, либо состоят из меньшего числа более крупных компонент. На основании этих признаков вычисляется простая эвристическая оценка, по которой блоку присваивается метка \texttt{text} или \texttt{non\_text}.

Полученные блоки сохраняются как отдельные изображения во временный каталог. При необходимости алгоритм формирует отладочную визуализацию, в которой найденные области отображаются поверх исходной страницы с указанием их типа и оценки уверенности. Это упрощает качественный анализ работы сегментатора и подбор параметров для различных классов документов.

Предложенный подход обладает рядом практических преимуществ. Во-первых, он вычислительно прост и не требует обучения на размеченной выборке. Во-вторых, метод хорошо согласуется с типичной вертикальной структурой научных и учебных изданий. В-третьих, он даёт интерпретируемый результат: каждая граница блока объясняется наличием достаточно широкой пустой полосы на странице.
Однако используемая схема сегментации основана на предположении о полноширинной блочной компоновке и потому не является универсальной. Она хуже подходит для многостолбцовой вёрстки, страниц со сложным обтеканием рисунков, боковыми примечаниями, плотными таблицами, а также для случаев, когда соседние элементы расположены очень близко друг к другу и не разделены выраженным светлым промежутком. Кроме того, классификация блоков на основе плотности и числа малых компонент является эвристической и потому не гарантирует идеального разделения текста и графики во всех случаях.

Тем не менее для рассматриваемой задачи такой метод оказывается практически полезным. Он позволяет быстро отделять крупные текстовые области от нетекстовых фрагментов страницы и тем самым направлять разные части документа в специализированные подсистемы последующей обработки. В контексте разрабатываемого программного комплекса сегментация по пустотам выступает промежуточным, но важным этапом, обеспечивающим согласование между анализом структуры страницы и дальнейшим распознаванием её содержимого.